\iffalse
For a vertex $v^t$, we define its \emph{neighborhood} as the set of previously arrived vertices that are co-clustered with $v^t$ in at least half of the input clusterings revealed up to time~$t$.
Formally,
\[
    \nbhd(v^t)
    \;=\;
    \left\{\, v^{t'} \;\middle|\; t' < t \;\land\; v^{t'} \sim_{\ge m/2} v^t \,\right\},
\]
where $v^{t'} \sim_{\ge m/2} v^t$ denotes that $v^{t'}$ and $v^t$ belong to the same cluster in at least $m/2$ of the input clusterings
$\m{C}_k^t$, for $k \in [m]$.
\fi

\iffalse
\paragraph{Algorithm Description.}
We are now ready to describe the $\oncon$ algorithm.
At time stamp~$t$, upon the arrival of vertex $v^t$, the algorithm proceeds as follows.

\begin{enumerate}
    \item If the neighborhood $\nbhd(v^t)$ intersects vertices from exactly one cluster in the current clustering $\m{C}^t$, then $v^t$ is assigned to that cluster.

    \item If the neighborhood $\nbhd(v^t)$ does not intersect any cluster in $\m{C}^t$, then $v^t$ is placed in a new singleton cluster.

    \item Otherwise, $v^t$ is assigned to a cluster in $\m{C}^t$ that maximizes the size of its intersection with $\nbhd(v^t)$.
    More specifically, $v^t$ is assigned to a cluster
    \[
        C_{\max} \in \arg\max_{C \in \m{C}^t} \; \card{C \cap \nbhd(v^t)}.
    \]
    Ties, if any, are broken arbitrarily.
\end{enumerate}
\fi

\iffalse

\begin{proof}
Fix a vertex $v^{t'}$ that arrived at time stamp $t' \le t$.
Suppose that upon arrival, the algorithm assigns $v^{t'}$ to a cluster
$C_k^{t'-1} \in \m{C}^{t'-1}$.

To bound the disagreement between $\m{C}^t$ and $\m{C}^{t^*}$,
we analyze pairs incident to $v^{t'}$.
Partition the set of good pairs $G$ incident to $v^{t'}$
into two subsets $D$ and $G \setminus D$, where $D$ consists of all
good pairs $(v^{t'},u)$ such that exactly one of the following holds:

\begin{itemize}
    \item $u \in \nbhd(v^{t'})$ but $u \notin C_k^{t'-1}$, or
    \item $u \notin \nbhd(v^{t'})$ but $u \in C_k^{t'-1}$.
\end{itemize}

Now consider any pair $(v^{t'}, w) \in G \setminus D$,
where $w$ arrived before $v^{t'}$.
By definition of $G \setminus D$, we have that either

\begin{itemize}
    \item $w \in \nbhd(v^{t'})$ and $w \in C_k^{t'-1}$, or
    \item $w \notin \nbhd(v^{t'})$ and $w \notin C_k^{t'-1}$.
\end{itemize}

In the first case, the algorithm places $v^{t'}$ together with $w$
in $\m{C}^{t'}$, and this relationship remains unchanged in $\m{C}^t$.
Since $w \in \nbhd(v^{t'})$, the pair $(v^{t'},w)$ is co-clustered
in at least $m/2$ input clusterings.
Because the pair is good, the optimal clustering $\m{C}^{t^*}$
also co-clusters $v^{t'}$ and $w$.
Thus, $\m{C}^t$ and $\m{C}^{t^*}$ agree on this pair.

In the second case, the algorithm separates $v^{t'}$ and $w$,
and since $w \notin \nbhd(v^{t'})$, the pair is separated
in at least $m/2$ input clusterings.
Again, because the pair is good, $\m{C}^{t^*}$ also separates them.
Thus, the two clusterings agree in this case as well.

Therefore, $\m{C}^t$ and $\m{C}^{t^*}$ agree on all pairs in
$G \setminus D$.
It follows that any disagreement between the two clusterings
must arise from either a bad pair or a pair in $D$.
Hence,
\[
    \dist(\m{C}^t, \m{C}^{t^*})
    \le
    |\overline{G}| + |D|.
\]

We now show that $|D| \le 2 |\overline{G}|$,
which completes the proof.
\stuck{I am stuck here. Previously, I thought this can be proved but now it seems impossible.}
\end{proof}
\fi
